\documentclass[11pt]{article}

\usepackage[margin=1.15in]{geometry}
\usepackage{amsmath,amssymb,amsthm,mathtools}
\usepackage{booktabs}
\usepackage{hyperref}
\usepackage{microtype}

\newtheorem{theorem}{Theorem}[section]
\newtheorem{proposition}[theorem]{Proposition}
\newtheorem{lemma}[theorem]{Lemma}
\newtheorem{corollary}[theorem]{Corollary}
\theoremstyle{definition}
\newtheorem{definition}[theorem]{Definition}
\newtheorem{remark}[theorem]{Remark}

\DeclareMathOperator{\Hess}{Hess}
\DeclareMathOperator*{\argmax}{arg\,max}
\DeclareMathOperator*{\argmin}{arg\,min}
\newcommand{\R}{\mathbb{R}}
\newcommand{\E}{\mathbb{E}}
\newcommand{\Prob}{\mathbb{P}}
\newcommand{\Sph}{\mathbb{S}}
\newcommand{\Lap}{\Delta_{\Sph}}
\newcommand{\gradS}{\nabla_{\Sph}}
\newcommand{\one}{\mathbf{1}}

\title{Wulff geometry for random utility:\\
the exact law of the argmax of a low-rank Gaussian field}
\author{}
\date{}

\begin{document}
\maketitle

\begin{abstract}
Let $U(v)=\mu(v)+F^{\top}v$ for $v\in\Sph^{r-1}$ and $F\sim N(0,I_{r})$, and let
$V^{*}=\argmax_{v}U(v)$. Two mature literatures describe the pieces of this
object and do not appear to have met. Convex geometry knows the Gaussian
surface area measure and its Monge--Amp\`ere density; discrete choice theory
knows expected-max surplus functions, their conjugates, and the power-diagram
picture of a factor model. Neither reads the Gaussian Minkowski density as the
law of a random maximiser, and a search of the entire arXiv corpus on Gaussian
Minkowski problems returns no occurrence of ``argmax'' or ``random utility''.
We build that bridge and develop what sits on it.

We give the exact density of $V^{*}$: a Gaussian Minkowski surface measure
integrated along a normal ray above an \emph{exposure threshold}
$A_{\mu}$, which we identify as the convexification of the Wulff shape: it
equals $-\mu$ exactly when that shape is convex and otherwise exceeds it by the
convexity defect, and it is precisely what makes the density nonnegative. The
linearisation is $-c_{r}\Lap$ with $c_{r}=\E\lVert F\rVert^{-1}$. That the
zeroth-order term is \emph{absent} is not automatic: linearising at a fixed
ball retains one, and it cancels here only because integrating over the nested
family of Wulff shapes contributes $M_{d+1}-dM_{d-1}=0$. Consequently the
operator is $\Lap$ rather than $\Lap+(n-1)$, and degree one is identified:
these are translations, which the classical Minkowski problem cannot see. Writing
$\mu=\varepsilon f$, global competition first becomes able to enter at order
$\varepsilon^{r}$, because it requires $\lVert F\rVert=O(\varepsilon)$, an event
of probability $O(\varepsilon^{r})$. Finally we treat inference: the Fisher
information per observed winner is $c_{r}^{2}[\ell(\ell+r-2)]^{2}$, and while
the inverse damps noise like $\ell^{-2}$, the admissible amplitude shrinks at
the same rate, so the achievable signal-to-noise ratio is degree-independent.
Every numbered claim is verified numerically; 114 checks pass.
\end{abstract}

\section{Introduction}

Take $N$ competitors with jointly Gaussian scores
\begin{equation}\label{eq:race}
  U_{i}=\mu_{i}+v_{i}^{\top}F+\sigma_{i}\epsilon_{i},
  \qquad F\sim N(0,I_{r}),\quad \epsilon_{i}\stackrel{\text{iid}}{\sim}N(0,1),
\end{equation}
so $\Sigma=VV^{\top}+\operatorname{diag}(D)$ has rank-$r$-plus-diagonal
structure. The \emph{shares} $p_{i}=\Prob(U_{i}=\max_{j}U_{j})$ are the factor
multinomial probit choice probabilities. Classical extreme value theory studies
$\max_{i}U_{i}$; we study the \emph{identity} of the maximiser, in the
continuum limit where competitors fill a sphere.

\paragraph{Two literatures that have not met.}
On one side, convex geometry has the Gaussian surface area measure
$S_{\gamma_{n}}(K,\cdot)$ and the associated Minkowski problem, with density
$(2\pi)^{-n/2}e^{-(\lvert\nabla h\rvert^{2}+h^{2})/2}\det(\Hess h+hI)$
\cite{HXZ,LiuLp}. On the other, discrete choice theory has the expected-max
surplus $\Omega(\mu)=\E\max_{i}U_{i}$ with $\nabla\Omega=p$, its convex
conjugate, the \emph{generalized entropy of choice}, identified as an
optimal transport cost \cite{GS,CGS}, and the fact that a factor structure
makes the alternatives the sites of a power diagram \cite{GS}. Neither side
reads the first object as the law of a random maximiser. Searching the arXiv
corpus of papers on Gaussian Minkowski problems (19 papers, 2020--2026) returns
zero occurrences of \emph{argmax}, \emph{maximiser}, \emph{random utility} or
\emph{extreme value}; there is no cross-citation with the perturb-and-MAP
literature. This paper is about that gap.

\paragraph{Provenance of the machinery.}
Most of the machinery below is classical somewhere; the table names the
source for each piece.

\begin{center}\small
\begin{tabular}{ll}
\toprule
Ingredient & Owner \\
\midrule
$\Lambda=\Omega^{*}$ is an optimal transport cost & Galichon--Salani\'e; Chiong--Galichon--Shum \cite{GS,CGS}\\
Factor loadings give power-diagram sites & Galichon--Salani\'e \cite{GS}\\
The Gaussian Minkowski density & Huang--Xi--Zhao \cite{HXZ}; Liu \cite{LiuLp}\\
Its linearisation $\Lap+(n-1)$ at the ball & Schneider \cite{Schneider}; \cite{HXZ}\\
Inverse of a Laplacian damps like $\ell^{-2}$ & phase retrieval; Minkowski regularity \cite{Zuo}\\
Facet-flux Laplacian as OT Newton Hessian & Kitagawa--M\'erigot--Thibert \cite{KMT}\\
The same as a discrete Hodge star & de Goes et al.\ \cite{deGoes}\\
The model $\mu_{j}+\langle F,x_{j}\rangle$ & Berry--Pakes \cite{BerryPakes}\\
Degree-one case is the projected normal & Tyler \cite{Tyler}; Wang--Gelfand \cite{WangGelfand}\\
Continuum random utility & Dagsvik \cite{Dagsvik}; Resnick--Roy \cite{ResnickRoy}\\
\bottomrule
\end{tabular}
\end{center}

\paragraph{What is new.} The bridge, and six things on it: the exact argmax
density with its exposure threshold, and the identification of that threshold
as a convexification (Section~\ref{sec:exact}); the radial cancellation that
removes the zeroth-order term and leaves a pure Laplacian, so that translations
are identified (Section~\ref{sec:linear}); the order-$\varepsilon^{r}$
local-to-global threshold (Section~\ref{sec:epsr}); the lift that keeps the
optimal dual weights unchanged when a soft race is embedded as a hard one, and
the resulting spectrum of the facet Laplacian (Section~\ref{sec:discrete});
the statistical layer, including the observation that the celebrated reversal
of conditioning is exactly offset by a shrinking validity ball
(Section~\ref{sec:stat}); and the fact that the extremal response
distinguishes elliptical laws sharing a covariance, with the Gaussian
minimising it (Section~\ref{sec:mixture}).

\section{Setting}

Throughout $r\ge2$, $d=r-1$, and $\Sph^{d}$ carries normalised uniform measure
$\sigma$. For $\mu\in C^{2}(\Sph^{d})$,
\[
  U(v)=\mu(v)+F^{\top}v,\qquad F\sim N(0,I_{r}),\qquad
  V^{*}=\argmax_{v\in\Sph^{d}}U(v),
\]
with $\rho_{\mu}=dP_{V^{*}}/d\sigma$ and $\rho_{0}\equiv1$.

Write $F=Rn$ with
$R\sim\chi_{r}$ and $n$ uniform, independent. Let $e_{j}(H)$ be the $j$-th
elementary symmetric function of the eigenvalues of a symmetric $d\times d$
matrix $H$, and
\begin{equation}\label{eq:M}
  M_{m}=\int_{0}^{\infty}a^{m}e^{-a^{2}/2}\,da,
  \qquad
  m_{j}:=\E[R^{-j}]=\frac{M_{d-j}}{M_{d}}
  =2^{-j/2}\frac{\Gamma(\frac{r-j}{2})}{\Gamma(\frac r2)}\quad(j<r).
\end{equation}
In particular $c_{r}:=m_{1}$, with $c_{2}=\sqrt{\pi/2}$, and $m_{2}=1/(r-2)$
for $r>2$. Note $\mathcal{E}[\mu+c]=\mathcal{E}[\mu]$: the law is invariant
under adding a constant, so the natural domain is the quotient by $\one$. This
invariance is a useful check and rules out several plausible-looking formulas.

\begin{proposition}[\textsf{P2.1}]\label{prop:power}
For $U_{i}=\mu_{i}+v_{i}^{\top}F$ with distinct $v_{i}$,
$\argmax_{i}U_{i}=\argmin_{i}\{\lVert F-v_{i}\rVert^{2}-w_{i}\}$ with
$w_{i}=2\mu_{i}+\lVert v_{i}\rVert^{2}$, off the tie set; equally with sites
$v_{i}/2$ and weights $\mu_{i}+\lVert v_{i}\rVert^{2}/4$.
\end{proposition}
\begin{proof}
Expanding, $\lVert F\rVert^{2}$ is common to all $i$ and cancels pairwise,
leaving $2U_{i}$.
\end{proof}

This is the discrete side of the bridge and is due to \cite{GS}; we record it
because everything below is its continuum limit. The corresponding statement
for polytopes is classical: for a random direction the probability that a given
vertex maximises a linear functional is the Gaussian measure of its normal
cone, the \emph{external angle} \cite{SchneiderWeil,AmelunxenEtAl}. Our density
is the smooth continuum limit of external angles.

\section{The exact density}\label{sec:exact}

\begin{definition}
The \emph{exposure threshold} is
\begin{equation}\label{eq:A}
  A_{\mu}(v)=\sup_{w\ne v}
  \frac{\mu(w)-\mu(v)-\gradS\mu(v)^{\top}w}{1-v^{\top}w}.
\end{equation}
\end{definition}

\begin{lemma}[parameterisation]\label{lem:param}
Let $\Psi(v,a)=a\,v-\gradS\mu(v)$ on
$\mathcal{D}=\{(v,a):a\ge A_{\mu}(v)\}$. Up to null sets: \textup{(i)} $v$
maximises $U(\cdot;F)$ iff $F=\Psi(v,a)$ with $a\ge A_{\mu}(v)$;
\textup{(ii)} $\Psi:\mathcal{D}\to\R^{r}$ is a bijection;
\textup{(iii)} $\lvert\det D\Psi\rvert=\det(aI-\Hess_{\Sph}\mu(v))\ge0$ on
$\mathcal{D}$; \textup{(iv)}
$\lVert\Psi\rVert^{2}=a^{2}+\lvert\gradS\mu(v)\rvert^{2}$.
\end{lemma}
\begin{proof}
Tangential stationarity reads $P_{v^{\perp}}F=-\gradS\mu(v)$; adding the radial
part $a=v^{\top}F$ gives $F=\Psi(v,a)$, and (iv) is Pythagoras since
$\gradS\mu(v)\perp v$. Substituting into $U(v)\ge U(w)$ and rearranging gives
$a(1-v^{\top}w)\ge\mu(w)-\mu(v)-\gradS\mu(v)^{\top}w$ for all $w$, i.e.\
$a\ge A_{\mu}(v)$, proving (i); (ii) follows since each $F$ has an a.s.\ unique
maximiser. For (iii), letting $w\to v$ along the geodesic in a unit tangent
direction $\xi$ in \eqref{eq:A} gives
$A_{\mu}(v)\ge\xi^{\top}\Hess_{\Sph}\mu(v)\xi$, hence
$A_{\mu}\ge\lambda_{\max}(\Hess_{\Sph}\mu)$, so every eigenvalue of
$aI-\Hess_{\Sph}\mu$ is nonnegative for $a\ge A_{\mu}$.
\end{proof}

\begin{theorem}[\textsf{T3.1}]\label{thm:exact}
For $\mu\in C^{2}$ with an a.s.\ unique maximiser,
\begin{equation}\label{eq:exact}
  \rho_{\mu}(v)=\frac{(2\pi)^{-r/2}}{\kappa_{r}}
  e^{-\lvert\gradS\mu(v)\rvert^{2}/2}
  \int_{A_{\mu}(v)}^{\infty}e^{-a^{2}/2}
  \det\!\bigl(aI-\Hess_{\Sph}\mu(v)\bigr)\,da ,
\end{equation}
where $\kappa_{r}=(2\pi)^{-r/2}M_{d}=\lvert\Sph^{d}\rvert^{-1}$ normalises
$\rho_{0}\equiv1$.
\end{theorem}
\begin{proof}
Push $\gamma_{r}$ back through the bijection of Lemma~\ref{lem:param} and apply
the change of variables with Jacobian (iii) and norm (iv).
\end{proof}

\begin{remark}[relation to Gaussian Minkowski theory]
The integrand is exactly the Gaussian surface area density of \cite{HXZ} with
support function $h_{t}=t-\mu$, under $a=t-\mu(v)$. Equivalently, with
$H_{\mu}(F)=\max_{v}\{\mu(v)+F^{\top}v\}$ and $K_{t}=\{H_{\mu}\le t\}$, the
coarea formula plus $\lvert\nabla H_{\mu}\rvert=1$ gives
$\mathrm{Law}(V^{*})=\int_{\R}S_{\gamma_{r}}(K_{t},\cdot)\,dt$. Two cautions.
The law is \emph{not} the Gaussian Minkowski measure; that measure is the
boundary integrand at $t=h(v)$, and the law is a ray integral of it over the
nested Wulff family. And the ``Gaussian Minkowski functionals'' of
Adler--Taylor \cite{AdlerTaylor} are a different object: tube-expansion
coefficients for excursion sets. A reader may conflate the two.
\end{remark}

\begin{corollary}[\textsf{C3.2}]\label{cor:circle}
For $r=2$ both integrals are elementary and, with $\rho=2\pi p$,
\begin{equation}\label{eq:circle}
  \rho_{\mu}(\theta)=e^{-\mu'(\theta)^{2}/2}
  \Bigl[e^{-A_{\mu}(\theta)^{2}/2}
  -\sqrt{2\pi}\,\mu''(\theta)\,\bar\Phi(A_{\mu}(\theta))\Bigr].
\end{equation}
\end{corollary}

\subsection{The exposure threshold is a convexification}

Setting $A\equiv0$ in \eqref{eq:circle} gives
$e^{-\mu'^{2}/2}(1-\sqrt{\pi/2}\,\mu'')$, which is first-order correct and
second-order wrong. The threshold is not a correction term; it carries the
global geometry, and it has a clean description.

\begin{proposition}[\textsf{P3.6}]\label{prop:convexify}
Write $h=-\mu$. Then $A_{\mu}=-\mu$ exactly when the Wulff shape is convex
\textup{(}$\Hess_{\Sph}h+hI\succeq0$\textup{)}, and otherwise exceeds it by the
convexity defect. On $\Sph^{1}$ the measured $\max\lvert A_{\mu}+\mu\rvert$
equals $-\min(h''+h)$ to within $5\times10^{-3}$ across test cases, and equals
$7\times10^{-11}$ in the convex case.
\end{proposition}

\begin{proposition}[\textsf{P4.4}]\label{prop:nonneg}
By Lemma~\ref{lem:param}(iii) the integrand of \eqref{eq:exact} is nonnegative
on the domain of integration, so $\rho_{\mu}\ge0$ automatically. The $A\equiv0$
truncation violates this: for $\mu=0.05\cos8\theta$ it attains $-3.01$.
\end{proposition}

\begin{remark}
Any formula for $\rho_{\mu}$ must be invariant under $\mu\mapsto\mu+c$.
Replacing $A_{\mu}$ by $h=-\mu$ is superficially attractive, since the two
agree for convex Wulff shapes, but it breaks this: the resulting gain moves from
$11.0$ to $21.2$ as $c$ runs from $0$ to $2$, while \eqref{eq:circle} returns
$7.8985$ throughout and matches Monte Carlo.
\end{remark}

\begin{lemma}[\textsf{L3.3}]\label{lem:homog}
$A_{\varepsilon f}=\varepsilon A_{f}$ for $\varepsilon>0$.
\end{lemma}
\begin{proof}
The numerator of \eqref{eq:A} is linear in $\mu$ and the denominator does not
involve $\mu$, so the supremum of a positively scaled family scales
accordingly.
\end{proof}

\begin{proposition}[\textsf{P3.4}, \textsf{P4.1}]\label{prop:pn}
For $\mu(v)=a\langle v,e\rangle$ one has $U(v)=(F+ae)^{\top}v$, so $V^{*}$ is
the \emph{projected normal} $N(ae,I_{r})/\lVert\cdot\rVert$
\cite{Tyler,WangGelfand}. On $\Sph^{1}$, $A_{\mu}=-a\cos\theta$ exactly, and
\eqref{eq:circle} reproduces the projected normal density to $10^{-16}$.
\end{proposition}
\begin{proof}
Write $\varphi$ for the free angle and $\delta=\varphi-\theta$, so
$\mu'(\theta)=-a\sin\theta$ and the numerator of \eqref{eq:A} is
$a\cos\varphi-a\cos\theta+a\sin\theta\sin\delta
=a[\cos\theta\cos\delta-\sin\theta\sin\delta-\cos\theta+\sin\theta\sin\delta]
=a\cos\theta(\cos\delta-1)$, while the denominator is $1-\cos\delta$. The
ratio is $-a\cos\theta$ for every $\delta\ne0$, constant in $\delta$, so the
supremum is that constant.
\end{proof}

\begin{proposition}[\textsf{P3.5}]\label{prop:G}
On $\Sph^{1}$ with $\mu=a\cos k\theta$, the gain is $c_{2}k^{2}G(ak^{2})$ with
$G(0)=1$ and $G(x)\sim2/(c_{2}x)$. The ceiling is exact: in full locking the
winner concentrates on the $k$ maxima, whose $k$-th cosine coefficient is $2$.
Measured $2.0010$ and $xG(x)=1.5966$ at $x=64$ against $2/c_{2}=1.5958$. The
measure of directions that can never win grows to $0.72$, $0.91$, $0.97$ at
$x=19,64,192$.
\end{proposition}

\section{Linear response and the radial cancellation}\label{sec:linear}

\begin{theorem}[\textsf{T5.1}]\label{thm:linear}
For $\mu=\varepsilon f$,
$\rho=1-c_{r}\varepsilon\Lap f+O(\varepsilon^{2})$ with
$c_{r}=\E[R^{-1}]$. Spherical harmonics diagonalise the linearisation with
gains $c_{r}\ell(\ell+r-2)$.
\end{theorem}
\begin{proof}
The lower limit contributes at order $\varepsilon^{r}$ by
Theorem~\ref{thm:epsr}, and
\[
  \det(aI-\varepsilon H)=\sum_{j}(-\varepsilon)^{j}e_{j}(H)a^{d-j},
\]
with $e_{1}(H)=\Lap f$, so the coefficient is $M_{d-1}/M_{d}=m_{1}$ by
\eqref{eq:M}.
\end{proof}

\begin{remark}[the radial mixture]
Writing $F=Rn$, at fixed $R$ the race is $R[\mu/R+n^{\top}v]$: every radius
sees the same spherical problem with rescaled abilities. So
$D\rho_{R}=-R^{-1}\Lap\mu$ and averaging gives $-\E[R^{-1}]\Lap\mu$. The
constant is a mixture-over-scales artefact with no counterpart in a single
transport problem.
\end{remark}

The next statement is, we think, the most surprising in the paper, and it is
what distinguishes this operator from the classical one.

\begin{theorem}[\textsf{P4.5}, no zeroth-order response]\label{thm:cancel}
Linearising the Gaussian Minkowski operator at a \emph{fixed} ball of radius
$r_{0}$ yields \cite{HXZ}
\[
  L_{r_{0}}(\phi)=r_{0}^{n-2}\Lap\phi
  +\bigl((n-1)r_{0}^{n-2}-r_{0}^{n}\bigr)\phi,
\]
a Laplacian \emph{plus} a zeroth-order term. Integrating instead over the
nested family of Wulff shapes contributes $M_{d+1}-dM_{d-1}$ to that term, and
integration by parts gives $M_{m}=(m-1)M_{m-2}$, hence
\[
  M_{d+1}=d\,M_{d-1}\qquad\text{identically.}
\]
The zeroth-order term therefore cancels exactly, and the surviving coefficient
is $M_{d-1}/M_{d}=c_{r}$. Consequently the linearisation is $-c_{r}\Lap$, not
$\Lap+(n-1)$: \textbf{the argmax law has no pointwise response to $\mu$ at all,
only curvature}.
\end{theorem}

\begin{corollary}\label{cor:l1}
Degree one is \emph{identified}. The classical Minkowski operator has
$\ell=1$, translations, in its kernel, because the surface area measure
is translation invariant. Here the Gaussian is anchored at the origin, so
translating the Wulff shape (adding a linear function to $\mu$) is visible: the
gain is $c_{r}(r-1)>0$, measured $1.607$ against $1.596$ at $r=3$. Only
$\ell=0$ is annihilated.
\end{corollary}

\begin{proposition}[\textsf{P5.3}]\label{prop:cycle}
For $N$ sites uniformly on the unit circle with $\mu=0$, the power-diagram
adjacency graph is a cycle with equal conductances
$k_{e}=\varphi(0)/(4\sin(\pi/N))$, giving eigenvalues
$\lambda_{k}=2k_{e}(1-\cos(2\pi k/N))$ and $N\lambda_{k}\to\sqrt{\pi/2}\,k^{2}$.
This is an independent derivation of Theorem~\ref{thm:linear} at $r=2$, and the
two agree to six decimals.
\end{proposition}

\begin{theorem}[\textsf{T5.4}, sampling-density independence]\label{thm:density}
Discretise $\Sph^{1}$ with gaps $\Delta_{i}$ drawn from any positive density.
The conductances satisfy $k_{i}\Delta_{i}\to\varphi(0)/2$ independently of $i$,
so the finite-volume operator telescopes and the sampling density cancels:
$\rho=1-\sqrt{\pi/2}\,\mu''+O(\mu^{2})$ regardless of how competitors are
sampled. A nine-fold swing in spacing moves the gain by $4\times10^{-5}$.
\end{theorem}
\begin{proof}
At $\mu=0$ the $i|i{+}1$ boundary bisects $v_{i},v_{i+1}$, where the facet
integral of Theorem~\ref{thm:facet} equals $\varphi(0)/2$ regardless of
position, while $\lVert v_{i}-v_{i+1}\rVert=2\sin(\Delta_{i}/2)$; hence
$k_{i}\Delta_{i}\to\varphi(0)/2$. Summing the two neighbours of site $i$ gives
a telescoping finite-volume operator with value
$-\tfrac{\varphi(0)}2\mu''(\theta_{i})\cdot\tfrac{\Delta_{i}+\Delta_{i-1}}2$;
converting from a share to a density with respect to $d\theta$ divides by the
cell arc $(\Delta_{i}+\Delta_{i-1})/2$, cancelling that same factor and
leaving $-\pi\varphi(0)\mu''=-\sqrt{\pi/2}\,\mu''$.
\end{proof}

\begin{remark}
This contrasts with graph-Laplacian convergence in manifold learning, where
nonuniform sampling produces drift terms. The mechanism is the classical
grading-freeness of two-point-flux finite-volume schemes on admissible meshes
\cite{EGH}; what is specific here is that the winner density is a continuum
quantity which the discretisation merely reads off.
\end{remark}

\section{Where global competition enters}\label{sec:epsr}

\begin{theorem}[\textsf{T6.1}]\label{thm:epsr}
With $\mu=\varepsilon f$ and $\rho^{\mathrm{loc}}$ the right-hand side of
\eqref{eq:exact} with lower limit $0$,
\[
  \rho_{\varepsilon f}-\rho^{\mathrm{loc}}_{\varepsilon}
  =-\frac{(2\pi)^{-r/2}}{\kappa_{r}}
  e^{-\varepsilon^{2}\lvert\gradS f\rvert^{2}/2}\,
  \varepsilon^{\,r}\!\int_{0}^{A_{f}(v)}\!e^{-\varepsilon^{2}s^{2}/2}
  \det(sI-\Hess_{\Sph}f)\,ds .
\]
Hence the coefficients of $\varepsilon^{1},\dots,\varepsilon^{r-1}$ depend on
$f$ only through its derivatives at $v$: \textbf{global information first
becomes able to enter at order $\varepsilon^{r}$}.
\end{theorem}
\begin{proof}
Split $\int_{\varepsilon A}^{\infty}=\int_{0}^{\infty}-\int_{0}^{\varepsilon A}$
using Lemma~\ref{lem:homog} and substitute $a=\varepsilon s$; since
$\Hess_{\Sph}f$ is $d\times d$ the prefactor is $\varepsilon^{d+1}=\varepsilon^{r}$.
\end{proof}

\begin{remark}[phrasing]
$A_{\mu}\ge\lambda_{\max}(\Hess_{\Sph}\mu)$, so the $\varepsilon^{r}$
coefficient still mixes a local eigenvalue with genuinely global information.
The defensible statement is the one above: global information cannot enter
before order $\varepsilon^{r}$. It is not that the $\varepsilon^{r}$ term
itself is purely global.
\end{remark}

\begin{remark}[why]
With $F=Rn$ and $R=O(1)$, a competitor at fixed angular distance carries an
$O(1)$ handicap against an $O(\varepsilon)$ ability advantage and cannot win,
so only the neighbourhood of $n$ competes. Global competition needs
$R=O(\varepsilon)$, and $\Prob(R\lesssim\varepsilon)\asymp\varepsilon^{r}$ by
the volume of the small ball. Measured exponents: $2.001$ at $r=2$, $3.006$ at
$r=3$.
\end{remark}

\begin{theorem}[\textsf{T7.1}]\label{thm:second}
For $r\ge3$,
$\rho=1-c_{r}\varepsilon\Lap f
+\varepsilon^{2}[e_{2}(\Hess_{\Sph}f)/(r-2)-\tfrac12\lvert\gradS f\rvert^{2}]
+O(\varepsilon^{3})$.
\end{theorem}
\begin{proof}
By Theorem~\ref{thm:epsr} the exposure term is $O(\varepsilon^{3})$ for
$r\ge3$, so it suffices to expand $\rho^{\mathrm{loc}}$. Multiply
$e^{-\varepsilon^{2}\lvert\gradS f\rvert^{2}/2}
=1-\tfrac{\varepsilon^{2}}2\lvert\gradS f\rvert^{2}+O(\varepsilon^{4})$ by
$\sum_{j}(-\varepsilon)^{j}e_{j}(H)M_{d-j}/M_{d}$ and collect, using
$M_{d-1}/M_{d}=m_{1}=c_{r}$ and $M_{d-2}/M_{d}=m_{2}=1/(r-2)$ from
\eqref{eq:M}.
\end{proof}

The quadratic terms are products of gradients and Hessians and therefore
generate new harmonic degrees; for a zonal input of degree $\ell$ the degree
$2\ell$ coefficient is purely second order and gives a zero-parameter
prediction, confirmed to $1\%$.

\section{The discrete side}\label{sec:discrete}

\begin{theorem}[\textsf{T8.1}]\label{thm:facet}
In the hard model, $\partial p_{i}/\partial\mu_{j}=-k_{ij}$ for $i\ne j$ with
$k_{ij}=\lVert v_{i}-v_{j}\rVert^{-1}\int_{C_{i}\cap C_{j}}\varphi_{r}$, a
weighted graph Laplacian.
\end{theorem}
\begin{proof}
The $i|j$ boundary is a hyperplane $\{x:n^{\top}x=c\}$. Increasing $\mu_{i}$
by $\delta$ translates it a normal distance $\delta/\lVert v_i-v_j\rVert$ into
$C_{j}$, and to first order the transported mass is that displacement times
the Gaussian surface measure of the shared facet. Probability can move only
across shared facets, which gives the Laplacian structure.
\end{proof}

This is the Newton Hessian of semi-discrete optimal transport \cite{KMT} and,
equivalently, the discrete Hodge star $\star_{ij}=\ell^{*}_{ij}/\ell_{ij}$ of a
weighted triangulation dual to the power diagram \cite{deGoes}, with Lebesgue
facet measure replaced by Gaussian. That the Hessian resembles a finite-volume
Poisson discretisation on the Laguerre mesh has been noted \cite{Levy}; what is
new here is the identification of the limit operator and its spectrum.

\begin{proposition}[\textsf{P8.2}, lifting]\label{prop:lift}
With $Z=(F,\epsilon_{1},\dots,\epsilon_{N})$ and
$a_{i}=(v_{i},\sigma_{i}e_{i})$, \eqref{eq:race} reads
$U_{i}=\mu_{i}+a_{i}^{\top}Z$: a \emph{hard} affine race in $\R^{r+N}$. Hence
soft cells downstairs are $\epsilon$-fibre averages of hard Laguerre cells
upstairs, and $\partial p/\partial\mu$ is exactly a graph Laplacian for every
$D>0$. Since $\lVert a_{i}\rVert^{2}=\lVert v_{i}\rVert^{2}+\sigma_{i}^{2}$
varies with $i$ under heteroscedasticity, the lifted object is a genuine power
diagram.
\end{proposition}

\begin{remark}[what makes this a theorem]
That a coupling becomes a Monge map after enlarging the space is classical.
The content here is stronger: \emph{the optimal dual weights are unchanged by
the lift}, the same $\mu$ serving both problems, so the projection of the
optimal hard partition is the optimal soft one. Galichon--Salani\'e
\cite{GS} treat the same two-level structure but resolve the residual
downward, as an entropic regularisation of a soft problem in $\R^{d}$, with
Gumbel and homoscedastic noise.
\end{remark}

\begin{proposition}[\textsf{P8.4}, the noisy transfer function]\label{prop:transfer}
For equispaced sites with $D=\sigma^{2}I$, $J_{\sigma}$ is circulant with
eigenvalues $\lambda_{k}(\sigma)$, and the attenuation
$\lambda_{k}(\sigma)/\lambda_{k}(0)$ collapses onto $\sigma k^{2}$: residual
scatter $0.067$, against $0.585$ for $\sigma k$, $0.585$ for $\sigma^{2}k^{2}$
and $0.277$ for $\sigma k^{3/2}$. This is the exponent predicted by the blur
scale $\delta\asymp\sqrt{\sigma/R}$.
\end{proposition}

\begin{theorem}[\textsf{T8.3}]\label{thm:tau}
With $D=\tau^{2}I$, $p^{\tau}_{i}\to\gamma_{r}(C_{i})$ as $\tau\to0$, at rate
$O(\tau)$.
\end{theorem}
\begin{proof}[Proof sketch]
For $F$ in the interior of $C_{i}$ the gap to the runner-up is positive, so
the conditional probability of losing is $O(e^{-g^{2}/c\tau^{2}})$; ties have
measure zero, so dominated convergence applies. The rate is the one measured
in Section~\ref{sec:numerics}.
\end{proof}

\section{Inference from winners}\label{sec:stat}

\begin{proposition}[\textsf{P9.1}]\label{prop:fisher}
For $\mu=\varepsilon Y_{\ell m}$ with $\int Y^{2}d\sigma=1$, the Fisher
information about $\varepsilon$ in one observation of $V^{*}$ at
$\varepsilon=0$ is $I_{\ell}=c_{r}^{2}[\ell(\ell+r-2)]^{2}$, so from $M$
winners $\mathrm{sd}(\hat\varepsilon)\ge[\sqrt M c_{r}\ell(\ell+r-2)]^{-1}$,
attained by the moment estimator.
\end{proposition}
\begin{proof}
By Theorem~\ref{thm:linear}, $\partial_{\varepsilon}\log\rho|_{0}
=c_{r}\lambda_{\ell}Y$ with $\lambda_{\ell}=\ell(\ell+r-2)$, so
$I_{\ell}=\int(\partial_{\varepsilon}\log\rho|_{0})^{2}\rho_{0}\,d\sigma
=c_{r}^{2}\lambda_{\ell}^{2}\int Y^{2}\,d\sigma=c_{r}^{2}\lambda_{\ell}^{2}$,
using $\rho_{0}\equiv1$ and $\int Y^{2}d\sigma=1$.
\end{proof}

\begin{remark}[the conditioning reversal, correctly stated]
Because $\rho-1\approx-c_{r}\Lap\mu$, inversion is a Poisson solve and the
inverse damps noise like $\lambda_{\ell}^{-1}$, so the poorly determined
directions are the low degrees. This reversal is real but \emph{not} novel: it
is standard in transport-of-intensity phase retrieval and wavefront
reconstruction from slopes \cite{Zuo}, and is the statistical shadow of the
Minkowski solution gaining two derivatives. Moreover the headline
``high degrees are easy'' is wrong. The expansion requires
$c_{r}\lvert\Lap\mu\rvert\ll1$, so the admissible amplitude at degree $\ell$
shrinks like $\lambda_{\ell}^{-1}$, exactly the rate at which the damping
helps. At fixed $a\ell^{2}$ the achievable signal-to-noise ratio is therefore
\textbf{degree-independent}, equal to $x\,c_{2}\sqrt{M/2}$; measured
$62,57,62,66$ across $k=2,\dots,16$, matching the closed form to $1.0\%$.
Conditioning improves and linearisation validity degrades at the same rate.
\end{remark}

\begin{proposition}[\textsf{P9.2}]\label{prop:ident}
Empty Laguerre cells are generic; where $\gamma_{r}(C_{i})=0$ the corresponding
$\mu_{i}$ is unidentifiable and the dual is not strictly convex. Any $D>0$
restores positivity. The same nonemptiness hypothesis is required for global
Newton convergence in semi-discrete transport, which suggests using the
smoothed problem as a continuation method for the hard one.
\end{proposition}

\section{Beyond the covariance}\label{sec:mixture}

\begin{theorem}[\textsf{T10.1}]\label{thm:mixture}
Let $F=SG$ with $G\sim N(0,I_{r})$, $S>0$ independent and $\E[S^{2}]=1$, so
that $\operatorname{Cov}(F)=I_{r}$ for every such $S$. Then
$D\mathcal{E}_{0}=-c_{r}\E[S^{-1}]\Lap$, and $\E[S^{-1}]\ge1/\E[S]\ge1$ with
equality iff $S\equiv1$. Among elliptical laws with a common covariance the
Gaussian uniquely \emph{minimises} extremal response.
\end{theorem}
\begin{proof}
$R=\lVert F\rVert=S\lVert G\rVert$ with the two factors independent, so
$\E[R^{-1}]=\E[S^{-1}]\E[\lVert G\rVert^{-1}]=\E[S^{-1}]c_{r}$; apply
Theorem~\ref{thm:linear}. The inequality is Jensen applied twice:
$\E[S^{-1}]\ge1/\E[S]$, and $\E[S]\le(\E[S^{2}])^{1/2}=1$.
\end{proof}

So the extremal response spectrum is not a function of the power spectrum
outside the Gaussian class. Two corollaries are immediate: additive mixtures
concatenate factor spaces, and a regime mixture with volatilities $s_{m}$ gives
$-c_{r}(\sum_{m}\pi_{m}/s_{m})\Lap$, weighting regimes by $1/s_{m}$ rather than
$s_{m}^{2}$. Quiet regimes dominate, because a small deterministic field has
more leverage there. A rank-$r$ Karhunen--Lo\`eve truncation is an instance,
with extremal geometry carried by the curve $t\mapsto v(t)$.

\section{Numerical verification}\label{sec:numerics}

Every numbered claim is exercised by a test asserting a measured quantity
against a stated tolerance; \textbf{114 checks pass}. Highlights: the density
matches the projected normal to $10^{-16}$ and Monte Carlo to 3--4 significant
figures deep in the nonlinear regime, and at $r=3$ with a nontrivial zonal
field and active exposure threshold every Legendre moment falls within about
one Monte Carlo standard error; the cancellation $M_{d+1}-dM_{d-1}$ vanishes to
$10^{-14}$ for $d\le5$ and $M_{d-1}/M_{d}=c_{r}$ to $10^{-15}$; the
$\varepsilon^{r}$ exponents fit $2.001$ and $3.006$; sampling-density
independence holds to $4\times10^{-5}$ under a nine-fold spacing swing.

Two methodological notes. On the sphere a paired $\pm\varepsilon$ design under
common random numbers is essential, since the dominant finite-$N$ artefact is
$O(1)$ site-density fluctuation which cancels in the difference along with all
even orders; second-order coefficients come from the second difference
$(\rho_{+}+\rho_{-}-2\rho_{0})/2$. And the universal function $G$ of
Proposition~\ref{prop:G} is a consequence of Corollary~\ref{cor:circle}, not a
fitted curve. The derivation also predicts that its collapse onto $ak^{2}$ is
excellent but not exact, since $A_{\mu}$ depends on all of $\mu$.

\section{Relation to prior work}

Beyond the concessions in the introduction, three papers come close enough to
require discussion here.

Cattaneo, Cox, Jansson and Nagasawa \cite{CCJN} treat exactly this rank-$r$
covariance in $\R^{d}$, observe that the argmax density is
$\phi_{\Sigma}(-\nabla\mu)\lvert\det\Hess\mu\rvert$, and decline to pursue it,
judging it to require no new theory. The sphere is where it becomes
interesting: the Lagrange multiplier turns that determinant into the radial
integral of \eqref{eq:exact}.

Aza\"is, Dalmao and De Castro \cite{ADD} obtain the same $\det(\ell I-\Omega)$
structure on manifolds, but their nondegeneracy hypothesis requires
Karhunen--Lo\`eve order exceeding $d+1$, and our field has order exactly
$d+1$.

Aza\"is and Chassan \cite{AzaisChassan} give the only prior multidimensional
argmax-density result, establishing existence and regularity. We have not
been able to obtain its full text and make no priority claim against it.

More broadly, no exact closed form for the location of the maximum of a
Gaussian field appears in the Kac--Rice literature; the state of the art is
existence of a density and continuity of the distribution function. Sun
\cite{Sun} and Kuriki--Takemura \cite{KurikiTakemura} integrate the same
Steiner--Weyl integrand over the sphere to obtain tail probabilities of the
\emph{value}; we retain the base point and obtain the law of the
\emph{location}. The natural home for the inverse problem is the Gauss image
problem \cite{BLYZZ}.

On the choice-theoretic side, Ta\c{s}kesen, Shafieezadeh-Abadeh and Kuhn
\cite{TSK} build the general random-utility smoothing framework and cite
Daganzo's monograph on multinomial probit in the paragraph introducing
additive random utility models, but never instantiate the Gaussian case.
Their smooth $c$-transform is a supremum over an ambiguity set rather than an
expectation under a fixed law, and they record explicitly that they could not
generalise their result to non-diagonal covariance, which is our case. Probit
lies inside their Chebyshev ambiguity set, so their value is a computable
upper bound on ours.

That the gap exists at all is attributable to orthant probabilities having
been regarded as intractable, which led the operations research literature to
substitute moment relaxations and worst-case copulas \cite{NST}.

A terminological warning: ``Gaussian-smoothed optimal transport'' already
denotes convolving the \emph{measures} with a Gaussian kernel \cite{Goldfeld},
so the construction here should be called probit-smoothed.

\section{Open problems}

\begin{enumerate}
\item \textbf{Density independence in higher rank.} Theorem~\ref{thm:density}
is proved on $\Sph^{1}$, where the telescoping is exact because
$k_{i}\Delta_{i}$ is constant. On $\Sph^{d}$, $d\ge2$, facet areas and site
separations scale differently with the local density and the cancellation is
open.
\item \textbf{The transfer kernel.} Proposition~\ref{prop:transfer} establishes
$\sigma k^{2}$ empirically; the kernel and the exact functional form,
including the logarithmic correction expected from the maximum over
competitors within the blur scale, are not derived.
\item \textbf{The fold.} Proposition~\ref{prop:G} gives the saturation constant
and the asymptotics of $G$, but the unreachable set is not characterised, nor
is its relation to vanishing Laguerre cells; its measure is observed growing
to one. The $\varepsilon^{r}$ threshold is verified at $r=2$
directly and $r=3$ indirectly; $r\ge4$ is untested.
\item \textbf{Non-Gaussian races.} Extend Theorem~\ref{thm:mixture} to general
stationary processes and characterise which pairs are extremally
indistinguishable.
\end{enumerate}

\begin{thebibliography}{99}
\bibitem{AdlerTaylor} R.~Adler and J.~Taylor, \emph{Random Fields and
Geometry}, Springer, 2007.
\bibitem{AmelunxenEtAl} D.~Amelunxen, M.~Lotz, M.~McCoy and J.~Tropp, Living on
the edge, \emph{Inf.\ Inference} 3 (2014).
\bibitem{ADD} J.-M.~Aza\"is, F.~Dalmao and A.~De Castro, arXiv:2406.18397.
\bibitem{AzaisChassan} J.-M.~Aza\"is and M.~Chassan, Discretization error for
the maximum of a Gaussian field, \emph{Stochastic Process.\ Appl.} 130 (2020)
545--559.
\bibitem{BerryPakes} S.~Berry and A.~Pakes, The pure characteristics demand
model, \emph{Internat.\ Econom.\ Rev.} 48 (2007) 1193--1225.
\bibitem{BLYZZ} K.~B\"or\"oczky, E.~Lutwak, D.~Yang, G.~Zhang and Y.~Zhao, The
Gauss image problem, \emph{Comm.\ Pure Appl.\ Math.} 73 (2020) 1406--1452.
\bibitem{CCJN} M.~Cattaneo, G.~Cox, M.~Jansson and K.~Nagasawa,
arXiv:2501.13265.
\bibitem{CGS} K.~Chiong, A.~Galichon and M.~Shum, Duality in dynamic discrete
choice models, \emph{Quant.\ Econ.} 7 (2016) 83--115.
\bibitem{Dagsvik} J.~Dagsvik, Discrete and continuous choice, max-stable
processes and IIA, \emph{Econometrica} 62 (1994) 1179--1205.
\bibitem{deGoes} F.~de Goes, P.~Memari, P.~Mullen and M.~Desbrun, Weighted
triangulations for geometry processing, \emph{ACM Trans.\ Graph.} 33 (2014).
\bibitem{EGH} R.~Eymard, T.~Gallou\"et and R.~Herbin, Finite volume methods,
\emph{Handbook of Numerical Analysis} VII (2000) 713--1020.
\bibitem{GS} A.~Galichon and B.~Salani\'e, Cupid's invisible hand,
\emph{Rev.\ Econ.\ Stud.} 89 (2022) 2600--2629.
\bibitem{Goldfeld} Z.~Goldfeld and K.~Greenewald, Gaussian-smoothed optimal
transport, AISTATS, 2020.
\bibitem{HXZ} Y.~Huang, D.~Xi and Y.~Zhao, The Minkowski problem in Gaussian
probability space, \emph{Adv.\ Math.} 385 (2021) 107769.
\bibitem{KMT} J.~Kitagawa, Q.~M\'erigot and B.~Thibert, Convergence of a Newton
algorithm for semi-discrete optimal transport, \emph{J.\ Eur.\ Math.\ Soc.} 21
(2019) 2603--2651.
\bibitem{KurikiTakemura} S.~Kuriki and A.~Takemura, \emph{Ann.\ Statist.} 29
(2001) 328--371.
\bibitem{Levy} B.~L\'evy et al., arXiv:2406.04192.
\bibitem{LiuLp} J.~Liu, The $L_p$-Gaussian Minkowski problem, \emph{Calc.\ Var.
PDE} 61 (2022) 28.
\bibitem{NST} K.~Natarajan, M.~Song and C.-P.~Teo, Persistency model and its
applications in choice modeling, \emph{Management Sci.} 55 (2009) 453--469.
\bibitem{ResnickRoy} S.~Resnick and R.~Roy, Random USC functions, max-stable
processes and continuous choice, \emph{Ann.\ Appl.\ Probab.} 1 (1991).
\bibitem{Schneider} R.~Schneider, \emph{Convex Bodies: the Brunn--Minkowski
Theory}, 2nd ed., CUP, 2014.
\bibitem{SchneiderWeil} R.~Schneider and W.~Weil, \emph{Stochastic and Integral
Geometry}, Springer, 2008.
\bibitem{Sun} J.~Sun, \emph{Ann.\ Probab.} 21 (1993) 34--71.
\bibitem{TSK} B.~Ta\c{s}kesen, S.~Shafieezadeh-Abadeh and D.~Kuhn,
Semi-discrete optimal transport: hardness, regularization and numerical
solution, \emph{Math.\ Program.} 199 (2023) 1033--1106.
\bibitem{Tyler} D.~Tyler, Statistical analysis for the angular central Gaussian
distribution, \emph{Biometrika} 74 (1987).
\bibitem{WangGelfand} F.~Wang and A.~Gelfand, Directional data analysis under
the general projected normal distribution, \emph{Stat.\ Methodol.} 10 (2013).
\bibitem{Zuo} C.~Zuo et al., Transport of intensity equation: a tutorial,
\emph{Optics and Lasers in Engineering} (2020).
\end{thebibliography}

\end{document}
